\section{الفصل~\ref{sec:muscles}. المخرج إلى العضلات}

بنية المخرج، أي الوحدات الحركية، وهامش أمان المشبك، والتشغيل بحسب الحجم، وزوج العضلات، وحلقة التمدد، ومولّدات الإيقاع، مقيسة مباشرة؛ وشرط استقرار الحلقة ذات التأخير مبرهنة؛ والنموذج الداخلي للجسم فرضية راسخة الأسس؛ وأرقام الأشكال حساب بنموذج الكتاب.

{\footnotesize
\setlength{\tabcolsep}{3pt}\renewcommand{\arraystretch}{1.15}
\begin{longtable}{>{\raggedright}p{0.5cm}>{\raggedright}p{4.0cm}>{\raggedright}p{5.6cm}>{\raggedright}p{2.3cm}>{\raggedright\arraybackslash}p{1.7cm}}
\toprule
م & القضية & التجربة وما قيس & المصدر & الحالة \\
\midrule\endfirsthead
\toprule
م & القضية & التجربة وما قيس & المصدر & الحالة \\
\midrule\endhead
1 & الوحدة الحركية: عصبون \nt{ACET} واحد يتحكم في مجموعة من الألياف، من بضع مئات إلى نحو ألفين؛ وفي العضلات ذات الضبط الدقيق تكون الوحدات أصغر & عضلات الإنسان، عدّ المحاور الحركية والألياف العضلية: في المتوسط نحو $340$ ليفًا لكل عصبون في العضلة بين العظام في اليد، ونحو $410$ في العضلة العضدية الكعبرية، ونحو $1930$ في عضلة الساق التوأمية الإنسية. ولا يذكر الفصل عددًا دقيقًا لأصغر الوحدات. & \cite{Feinstein1955mu} & مُقاس \\
2 & المشبك على العضلة يطلق من الناقل العصبي أضعاف ما يلزم لإطلاق الليف & مراجعة للقياسات على المشابك العصبية العضلية: هامش الأمان (نسبة الناقل المُطلَق إلى العتبي) لدى الثدييات البالغة عادة $3\text{--}5$؛ ولدى الإنسان يقوم الهامش على طيات غشاء المتلقي أكثر مما يقوم على حجم الإطلاق. & \cite{WoodSlater2001} & مُقاس \\
3 & النفضة~\eqref{eq:twitch} تتصاعد خلال زمن $T_c$ من رتبة $50\ms$ & وحدات حركية مفردة في يد الإنسان عند جهد إرادي، والقوة مُتوسَّطة على نبضات الوحدة: زمن التصاعد $30\text{--}100\ms$، وهو أقل من $70\ms$ لدى $80\,\%$ من الوحدات. والدالة $(t/T_c)e^{1-t/T_c}$ تقريب من نموذج لمجمع الوحدات. & \cite{MilnerBrown1973rec, Fuglevand1993} & مُقاس \\
4 & مجموع النفضات~\eqref{eq:forcesum}: عند $5$ و$15$ و$40$ نبضة في الثانية تكون القوة $0.2\text{--}1.1F_1$ و$1.8\text{--}2.2F_1$ ونحو $5.4F_1$ (الشكل~\ref{fig:tetanus}) & حساب بـ\foreignlanguage{english}{\texttt{models/\allowbreak muscle\_\allowbreak models.py}}، $T_c=50\ms$: $0.21\text{--}1.10$ و$1.76\text{--}2.19$ و$5.28\text{--}5.49\,F_1$ في آخر $0.3\,\text{s}$. والجمع الخطي تبسيط: لا تشبّع للعضلة الحقيقية في النموذج. & \foreignlanguage{english}{\texttt{models/\allowbreak muscle\_\allowbreak models.py}} & نموذج \\
5 & تُشغَّل الوحدات بحسب الحجم؛ وأضعفها أضعف من أقواها مئة مرة & الجذور البطنية لدى القط: مع دخل منعكس متزايد تفرغ أولًا العصبونات ذات النبضات الصغيرة في الجذر، أي الصغيرة. العضلة بين العظام في يد الإنسان: قوة نفضة الوحدات من $0.1$ إلى $10\,\text{g}$، وتزداد شبه خطيًا مع القوة التي تُشغَّل عندها الوحدة. & \cite{Henneman1957, MilnerBrown1973rec} & مُقاس \\
6 & خطوة القوة تتناسب مع القوة، $\Delta F/F\approx q-1$~\eqref{eq:sizeprinciple}: فاصلة عائمة & مستنتج في الفصل من المتوالية الهندسية للقوى. ويتفق مع التجربة: عند الجهود الكبيرة يُشغَّل لزيادة القوة نفسها عدد أقل بكثير من الوحدات الجديدة. & استنتاج في الفصل؛ \cite{MilnerBrown1973rec} & نظرية \\
7 & انتشار القوة يزداد بما يتناسب مع القوة نفسها & جهد إرادي متساوي القياس لعضلة الإبهام: الانحراف المعياري للقوة يزداد خطيًا مع متوسط القوة؛ وعند الجهد نفسه المستثار بالتنبيه الكهربائي يبقى الانتشار ثابتًا. نموذج المجمع: الازدياد الخطي يعطيه تشغيل الوحدات بترتيب القوة. & \cite{Jones2002sdn} & مُقاس \\
8 & نعومة الحركات نتيجة للضجيج المعتمد على الإشارة & نموذج: المسار الذي يقلل انتشار النقطة النهائية تحت هذا الضجيج يعيد إنتاج المسارات المقيسة للعين واليد. ولا توجد تجربة مباشرة تميّز هذا السبب من غيره. & \cite{HarrisWolpert1998} & فرضية \\
9 & الفرق بين قوتي زوج العضلات يحدد العزم، ومجموعهما الصلابة~\eqref{eq:antagonists} & نظرية التحكم في الممانعة بالشد المشترك. ذراع الإنسان: صلابة كل مفصل، مقيسة بدفعات صغيرة، تتناسب تقريبًا مع العزم فيه؛ والنسب المختلفة للشد المشترك تغيّر شكل قطع الصلابة الناقص لليد واتجاهه. & \cite{Hogan1984, GomiOsu1998stiff} & مُقاس \\
10 & مستشعر التمدد يثير عصبون \nt{ACET} الخاص به، ويطفئ المضادة عبر وسيط مثبِّط (الشكل~\ref{fig:antagonists}) & الحبل الشوكي لدى القط: وسيط واحد تثيره ألياف مستشعرات التمدد يعطي كمونات تثبيطية أحادية المشبك في عصبونات \nt{ACET} للعضلة المضادة، بتأخير مشبكي $0.28\text{--}0.42\ms$. ولم يُحدَّد في هذه التجربة الناقلُ العصبي للوسيط (الغليسين). & \cite{Jankowska1972Ia} & مُقاس \\
11 & تأخير الحلقة: الشوكية $25\text{--}30\ms$، وعبر القشرة أطول بمرة ونصف إلى ثلاث مرات، وعبر العين $100\text{--}200\ms$ & ذراع الإنسان، دفعة للمفصل: استجابة عضلية قصيرة بعد $20\text{--}45\ms$، وطويلة بعد $45\text{--}105\ms$، وإرادية بعد $120\text{--}180\ms$. إزاحة الموضع المرئي للإصبع أثناء الحركة: التصحيح بعد $160\ms$ في المتوسط. & \cite{Pruszynski2008rapid, SaundersKnill2003vis} & مُقاس \\
12 & $\dd x/\dd t=-Gx(t-d)$ مستقر عند $Gd<\pi/2$~\eqref{eq:delaystable}؛ وعلى الحد التردد $1/(4d)$ & مبرهنة عن جذور معادلة ذات تأخير. تحقق بـ\foreignlanguage{english}{\texttt{models/\allowbreak muscle\_\allowbreak models.py}}، $d=30\ms$: عند $3\,\text{s}$ السعة $3\cdot10^{-6}$ عند $G=40$، و$0.13$ عند $G=50$، و$38$ عند $G=55$ في الثانية (الحد $52$). & \cite{Hayes1950} & نظرية \\
13 & الرعشة الفسيولوجية $8\text{--}12\,\text{Hz}$؛ وحلقة التمدد أحد مصادرها & مراجعة للقياسات: الرعشة الدقيقة في نطاق نحو $10\,\text{Hz}$ موجودة لدى الأصحاء؛ ومنشؤها مختلط: تذبذبات مركزية، وخصائص تفريغ الوحدات، والرنين الميكانيكي، ورنين الحلقة المنعكسة، ويغلب أحدها أو غيره بحسب الظروف. & \cite{McAuleyMarsden2000} & فرضية \\
14 & يتنبأ الدماغ بنتيجة الأمر بالنموذج الداخلي للجسم & يمسك الإنسان ثقلًا بالقرص ويحرك ذراعه: مع الحمل القصوري واللزج والمرن تتغير قوة القبض متزامنة مع الحمل، لا بتأخير الحلقة، أي أن الحمل متنبَّأ به. وتجمع المراجعة تجارب كهذه؛ ولا ملاحظة مباشرة للنموذج نفسه في العصبونات. & \cite{FlanaganWing1997grip, WolpertGhahramani2000} & فرضية \\
15 & إيقاع المشي والتنفس والمضغ تولّده شبكات محلية عند دخل ثابت & مراجعة للتجارب: الجهاز العصبي المعزول (الحبل الشوكي، العقد) يُصدر نمطًا حركيًا إيقاعيًا بلا دخل إيقاعي وبلا تغذية راجعة من العضلات. & \cite{MarderBucher2001} & مُقاس \\
16 & التثبيط المتبادل مع التعب~\eqref{eq:halfcentre} يعطي إيقاعًا بدور $0.84\,\text{s}\approx2\tau_a$ (الشكل~\ref{fig:halfcentre}) & مبرهنة: الشبكة ذات التثبيط المتبادل والتكيّف التي لا حالة مستقرة لها عند دخل ثابت تتذبذب حتمًا. حساب بـ\foreignlanguage{english}{\texttt{models/\allowbreak muscle\_\allowbreak models.py}} ($I=1$، $\beta=2$، $b=2$، $\tau=20\ms$، $\tau_a=0.4\,\text{s}$): الدور $0.84\,\text{s}$. أما أن المولّدات الحقيقية مبنية هكذا بالذات فلم يُبيَّن. & \cite{Matsuoka1985osc} & نموذج \\
\bottomrule
\end{longtable}}
